% GENERATED FILE. Edit canonical lessons, then run npm run generate:books.
% canonical-source-hash: fnv1a64:e5e120e17bcb588f
% canonical-lessons: ML-C19-vayassu

\chapter{Asking Someone's Age}
\label{ch:ml-asking-age}

This chapter is generated from the canonical micro-lessons used by Language Ladder.
Each section stays independently resumable and preserves the authored prerequisite order.

\section[\ml{നിങ്ങൾക്ക്} \ml{എത്ര} \ml{വയസ്സുണ്ട്}?]{\ml{നിങ്ങൾക്ക്} \ml{എത്ര} \ml{വയസ്സുണ്ട്}? — "to you, how much age exists?"}

\label{lesson:ML-C19-vayassu}



\begin{quote}
\textbf{Warm-up.} {[}PAUSE 2s{]} Same Sanskrit word as Kannada and Telugu — but Malayalam builds a genuinely different sentence around it.
\end{quote}

\subsection*{You'll want to know}
\textbf{\ml{വയസ്സ്}} (\emph{vayass}) = "\textbf{age}" — the same Sanskrit \textbf{\dv{वयस्}} (\emph{vayas}, "age, vigor," PIE cousin of Latin \textbf{vīs}) behind Kannada's \emph{vayassu} and Telugu's \emph{vayasu}.

\begin{grammarlens}[title={Grammar Lens: dative subject + "exists"}]
\begin{quote}
\textbf{\ml{നിങ്ങൾക്ക്} \ml{എത്ര} \ml{വയസ്സുണ്ട്}?} — "How old are you?" (formal/plural) —
\end{quote}

literally "\textbf{to you, how much age exists}?" (\emph{ningaḷkku ethra vayassuṇṭŭ?} — from \emph{ningaḷkku}, "to you" {[}dative{]} + \emph{ethra vayassu}, "how much age" + \emph{uṇṭŭ}, "exists.")

\begin{quote}
\textbf{\ml{എനിക്ക്} \ml{ഇരുപത്} \ml{വയസ്സുണ്ട്}.} — "I am twenty years old." — literally "\textbf{to}
\end{quote}

me twenty age exists\textbf{." (\emph{enikku irupatu vayassuṇṭŭ.})}

This is a real grammatical departure from Kannada and Telugu's plain "your age how-much?" Malayalam puts the \textbf{person in the dative case} ("to you," not "your") and adds \textbf{\ml{ഉണ്ട്}} (\emph{uṇṭŭ}, "exists, there is") — the same existential verb Malayalam uses for ordinary possession and existence sentences generally. So the literal shape is "\textbf{to me, twenty years' worth of age exists}" — a dative-subject-plus-existential-verb construction, genuinely different from its two Dravidian cousins even though the age-word itself is identical.
\end{grammarlens}

\subsection*{Guided Practice}
{[}PAUSE 1s{]}

\begin{itemize}
\raggedright
  \item {[}YOU SAY: "vayassu" — age{]}
  \item {[}YOU SAY: "ningaḷkku ethra vayassuṇṭŭ?" — how old are you?{]}
  \item {[}YOU SAY: "enikku irupatu vayassuṇṭŭ" — I am twenty, "to me twenty age exists"{]}
\end{itemize}

\begin{tcolorbox}[breakable,colback=teal!4,colframe=teal!35!black,title={Wrap-up recall}]
{[}PAUSE 3s{]} Is Malayalam's \textbf{\ml{വയസ്സ്}} the same Sanskrit word as Kannada's and Telugu's? (\textbf{Yes}.) Does Malayalam ask age the same simple way Kannada and Telugu do? (\textbf{No} — it uses a \textbf{dative} subject, "to you," plus the existential verb \textbf{\ml{ഉണ്ട്}}, "exists.") What does \emph{enikku irupatu vayassuṇṭŭ} literally mean? ("\textbf{To me twenty age exists}.")
\end{tcolorbox}
